% SOURCE_STATEMENT: "Does every subfactor arise in conformal field theory?"
% Preparation and final self-audit: CONTRIBUTING.md prompts C and D; source/audit artifacts retained in the local physics_statements_2026_10_11 batch.
% !TeX program = lualatex
\documentclass[11pt]{article}
\usepackage[margin=25mm]{geometry}
\usepackage{fontspec}
\setmainfont{Latin Modern Roman}
\usepackage{amsmath,amssymb,mathtools}
\usepackage{unicode-math}
\setmathfont{Latin Modern Math}
\usepackage{xurl,hyperref}
\hypersetup{colorlinks=true,urlcolor=blue}
\setcounter{secnumdepth}{0}
\setlength{\parindent}{0pt}
\setlength{\parskip}{5pt}
\setlength{\emergencystretch}{3em}
\newcommand{\sref}[2]{\hyperlink{source:#1:#2}{[S#2]}}
\newcommand{\cataloguescope}[1]{\paragraph{Recorded target.}#1}
\newcommand{\R}{\mathbb R}
\newcommand{\C}{\mathbb C}
\newcommand{\N}{\mathbb N}
\newcommand{\Z}{\mathbb Z}
\newcommand{\Q}{\mathbb Q}
\newcommand{\F}{\mathbb F}
\newcommand{\eps}{\varepsilon}
\begin{document}
\section{Conformal Field Theory Realizations of Subfactors}\label{statement}

\subsection{Definitions and mathematical statement}

\paragraph{Shared notation.}$\mathbb N=\{1,2,\ldots\}$; the standard number fields have their usual meanings. All additional conventions are stated locally.


\paragraph{Definitions and conventions.}
A subfactor is an inclusion $N\subset M$ of von Neumann algebras with trivial centers. Its index and standard invariant encode the size and higher relative-commutant structure of the inclusion. In an operator-algebraic conformal field theory, interval or local observable algebras and sectors can produce subfactor inclusions. To say a subfactor arises in conformal field theory requires identifying the inclusion or invariant that the construction realizes. The primary report concerns finite Jones index, with the intended conformal realization recovering the standard invariant or at least its quantum double. This finite-index context is explicit in Vaughan Jones\textquotesingle s contribution, even though the catalogue sentence abbreviates it.

\paragraph{Target retained from the canonical source catalogue.}
\paragraph{Catalogue wording retained for traceability.}
Does every subfactor arise in conformal field theory?

\paragraph{Source scope and qualification.}
The primary contribution (PDF page 5, printed page 853) supplies the finite-index and realization qualifications omitted from the abbreviated catalogue sentence. ID 30004278 revisits the same broad construction programme, and ID 30002840 is the Asaeda--Haagerup instance; equivalence of every realization convention has not been established.
\subsection{Short English statement}
Can every subfactor be realized by a conformal field theory?
\subsection{Sources}
\begin{itemize}
\item[\hypertarget{source:statement:1}{S1}] ULAM.AI and the UnsolvedMath Contributors, UnsolvedMath, catalogue entry. Catalogue curation under CC BY 4.0; definitions and English exposition added locally.
\url{https://huggingface.co/datasets/ulamai/UnsolvedMath}.
\item[\hypertarget{source:statement:2}{S2}] Subfactors and Conformal Field Theory (2015), Oberwolfach Reports, DOI 10.4171/owr/2015/16, pp. 849--926. Primary report accessed 2026-10-11; relevant PDF pages 5.
\url{https://doi.org/10.4171/owr/2015/16}.
\item[\hypertarget{source:statement:3}{S3}] Bischoff--Davydov--Kawahigashi--Longo, index below four (catalogue-listed context; not independently consulted in this preparation).
\url{https://arxiv.org/abs/1506.02606}.
\end{itemize}
\label{end:statement}
\end{document}
